\chapter{Collections}

\begin{figure}[htbp]
    \centering
    \includegraphics[width=\textwidth]{JCollFramewk.pdf} 
    \caption{Gerarchia delle interfacce del Java Collection Framework}
    \label{fig:collection-hierarchy}
\end{figure}

\section{Separazione tra Interfaccia e Implementazione}

Il \textit{Java Collections Framework} adotta un approccio fondamentale della programmazione a oggetti: la netta separazione tra le \textbf{interfacce} (che definiscono il comportamento) e le \textbf{implementazioni} concrete (che definiscono le strutture dati e le prestazioni). Per illustrare bene questo concetto e i suoi vantaggi, esaminiamo il seguente esempio. 

Una \texttt{Queue} (coda) specifica un comportamento \textbf{FIFO} (\textit{First In, First Out}): aggiunta in coda, rimozione in testa e monitoraggio della dimensione.

\begin{lstlisting}[caption={Interfaccia minimale di una Queue}]
public interface Queue<E> {
    void add(E e);
    E remove();
    int size();
}
\end{lstlisting}

L'interfaccia non dice nulla su come i dati siano memorizzati. Esistono due implementazioni principali nel JDK:
\begin{enumerate}
    \item \textbf{\texttt{ArrayDeque}}: Utilizza un \textit{circular array} (array circolare, vedi in appendice \ref{appendix circ array}).
    \item \textbf{\texttt{LinkedList}}: Utilizza una lista concatenata di nodi.
\end{enumerate}

Il vantaggio principale di questa separazione è la flessibilità. Nel codice backend, è considerata una \textit{best practice} utilizzare il tipo interfaccia per dichiarare le variabili, limitando l'uso della classe concreta solo al momento della costruzione.

\begin{lstlisting}[caption={Utilizzo corretto delle interfacce}]
// Uso l'interfaccia come tipo della variabile
Queue<Customer> expressLane = new ArrayDeque<>(100);
expressLane.add(new Customer("Harry"));

// Se cambio idea, devo cambiare solo la riga del costruttore
Queue<Customer> expressLane = new LinkedList<>();
\end{lstlisting}

\begin{bestpractice}[Coding to an Interface]
Usare il tipo dell'interfaccia (es. \texttt{Queue}) invece del tipo concreto (es. \texttt{ArrayDeque}) protegge il codice da dipendenze eccessive. Se accidentalmente si utilizzasse un metodo specifico di \texttt{ArrayDeque} non presente in \texttt{Queue}, il compilatore segnalerebbe l'errore, garantendo che il codice rimanga facilmente intercambiabile.

In Java 6, quando fu introdotta \texttt{ArrayDeque}, i programmatori che avevano "programmato per interfacce" poterono migliorare istantaneamente le prestazioni delle loro applicazioni cambiando una singola riga di codice (il costruttore), lasciando intatto tutto il resto della logica aziendale.
\end{bestpractice}

Perché scegliere l'una o l'altra? L'interfaccia è agnostica riguardo alle prestazioni, ma le implementazioni no:

\begin{itemize}
    \item \textbf{\texttt{ArrayDeque}}: È generalmente più performante. Ha un impatto minore sul \textit{Garbage Collector} perché non deve creare un oggetto "Nodo" per ogni elemento inserito. Tuttavia, richiede la riallocazione dell'array quando la capacità finisce.
    \item \textbf{\texttt{LinkedList}}: Non ha problemi di capacità finita, ma ogni elemento aggiunto comporta la creazione di un nuovo oggetto in memoria, rendendola meno efficiente in scenari ad alto volume di dati.
\end{itemize}



